\originalchapter{多项式恒等式与素性检验}
\originalsection{Fermat 检验的局限}
素数 $p$ 的非零剩余类成阶 $p-1$ 群, 因而 $a^{p-1}\equiv1\pmod p$ 对 $p\nmid a$ 成立. 若某个与 $n$ 互素的 $a$ 不满足该同余, 则 $n$ 合成; 反向不成立. $561=3\cdot11\cdot17$ 对任意整数 $a$ 满足 $a^{561}\equiv a\pmod{561}$: 在每个素因子模下, 若 $a$ 为0结论成立, 否则群阶2,10,16均整除560, 再用中国剩余定理. 这种合数称 Carmichael 数.
\originalsection{多项式素性判别}
\begin{proposition}
设 $n\ge2$ 且 $\gcd(a,n)=1$. 则 $n$ 为素数当且仅当 $(x+a)^n=x^n+a$ 在 $(\Z/n\Z)[x]$ 中成立.
\end{proposition}
\begin{proof}
素数情形由 Frobenius 恒等式及 $a^n=a$. 合数情形取素数 $p\mid n$, 写 $n=p^km$, $p\nmid m$. $x^p$ 的系数为 $\binom np a^{n-p}$. 公式
\[
\binom np=\frac np\binom{n-1}{p-1}
\]
中第二因子模 $p$ 为 $\binom{-1}{p-1}\equiv1$, 因逐项分子为 $-1,-2,\ldots,-(p-1)$ 且分母可逆. 因而该二项系数恰含 $p^{k-1}$, 不被 $p^k$ 整除, $a$ 又不含 $p$, 所以该系数模 $n$ 非零. $p<n$, 右侧无 $x^p$ 项, 恒等式失败.
\end{proof}
直接算此恒等式需要处理 $n$ 个系数, 并未得到关于输入位数的快速算法. AKS 通过模 $x^r-1$ 压缩计算, 其完整参数界与正确性证明不在本册覆盖范围. 本册不复述原讲义未经此处证明的复杂度指数.
\originalsection{平方根见证}
\begin{proposition}
若奇数 $n$ 模下存在 $x\not\equiv\pm1$ 而 $x^2\equiv1$, 则 $n$ 合成, 并可从 $\gcd(x-1,n)$ 得到真因子.
\end{proposition}
\begin{proof}
素数域中 $(x-1)(x+1)=0$ 强迫 $x=\pm1$. 现设 $d=\gcd(x-1,n)$. 若 $d=1$, 则 $x-1$ 可逆, 推得 $x+1=0$ 模 $n$, 矛盾; 若 $d=n$, 则 $x=1$, 也矛盾. 所以 $1<d<n$.
\end{proof}
必须注意反命题错误: 模9的1只有两个平方根1,8, 但9是合数. 一般奇素数幂 $p^k$ 的1也只有两根, 因 $\gcd(x-1,x+1)$ 只可能含因子2, 故整个 $p^k$ 必整除其中一项. 只有含至少两个不同奇素因子的 $n$, 中国剩余定理才给至少四个符号组合根. 原18.703第22讲相关引理的“合数必至少四根”须加此条件.

设 $n\ge3$ 为奇数, 写 $n-1=2^sd$, $d$ 奇, $s\ge1$. 对与 $n$ 互素的底数 $a$, 若 $n$ 素, 序列 $a^d,a^{2d},\ldots,a^{2^sd}=1$ 必或者第一项已为1, 或在首次出现1的前一项为 $-1$. 这是域内平方根仅 $\pm1$ 的结果. 因而若 $a^d\not\equiv1$ 且所有前面平方链项都不为 $-1$, 可判 $n$ 合成. 这是强伪素数检验的确定见证方向; 未证明的错误概率上界不在此宣称.
\originalsection{带解答练习}
\begin{exercise}用平方根见证分解15, 并解释为什么同一办法不能仅靠模9的平方根判合数.\end{exercise}
\begin{solution}
$4^2\equiv1$ 模15但4不是 $\pm1$, $\gcd(4-1,15)=3$ 给真因子. 模9仅1和8满足平方为1, 因而没有非平凡平方根见证; 这不妨碍试除3或其他底数的幂检验判合数.
\end{solution}
